% Partie : théorie de la démonstration
% Chapitre : élimination des coupures
% Section : interpolation

\documentclass[../../../include/open-logic-section]{subfiles}

\begin{document}

\olfileid[fr]{pt}{cut}{itp}

\olsection{Lemme de Maehara et théorème d'interpolation de Craig}

Le théorème d'interpolation, dû à William Craig, affirme que, si
$!A \Entails !B$, il existe !!a{sentence}~$!C$, appelé
\emph{interpolant}, telle que $!A \Entails !C$ et $!C \Entails !B$.
De plus, on peut choisir l'interpolant~$!C$ de sorte que toutes les
!!{constant}s et tous les !!{predicate}s qui y apparaissent figurent à la fois
dans~$!A$ et dans~$!B$. Nous supposons qu'aucun !!{function} n'intervient.
Sans cette restriction sur le vocabulaire, $!A$ et $!B$ seraient
eux-mêmes des interpolants.

Le théorème d'interpolation est valable en logique classique du premier
ordre. On l'a qualifié de « dernière propriété importante de la logique
du premier ordre » à avoir été découverte. Il a des applications
importantes en théorie des modèles et en raisonnement automatisé.
Sa motivation et son application initiales relevaient de la philosophie
des sciences : il permet d'étudier au moyen de quelles notions primitives
une théorie peut, ou ne peut pas, être axiomatisée. Il implique aussi
le théorème de définissabilité de Beth : lorsqu'une théorie peut définir
implicitement une notion, elle peut également la définir explicitement.

Comme beaucoup de résultats de logique mathématique, le théorème
d'interpolation se démontre aussi bien par les méthodes de la théorie
des modèles que par celles de la théorie de la démonstration. Ces
dernières offrent un avantage : elles établissent l'existence d'un
interpolant et fournissent un algorithme qui le calcule à partir d'une
!!{proof} de $!A \Entails !B$, par exemple d'une !!{proof} dans
\Log{G3c} de $!A \Sequent !B$.

Notons $\Lang L(!A)$ l'ensemble des !!{constant}s et des !!{predicate}s
de~$!A$, et posons $\Lang L(\Gamma) = \bigcup \Setabs{\Lang L(!A)}{!A \in
\Gamma}$. Dans toute cette section, les !!{formula}s ne contiennent
aucun !!{function}.

\begin{thm}[Théorème d'interpolation de Craig]\ollabel{thm:interpolation}
Soient $\Lang L(!A)$ le !!{language} de~$!A$ et $\Lang L(!B)$ celui
de~$!B$. Si $\Log{G3c} \Proves !A \Sequent !B$, il existe une
!!{formula}~$!C$ dont le !!{language} vérifie
$\Lang L(!C) \subseteq \Lang L(!A) \cap \Lang L(!B)$, telle que
$\Log{G3c} \Proves !A \Sequent !C$ et
$\Log{G3c} \Proves !C \Sequent !B$.
\end{thm}

Pour exploiter la structure des !!{proof}s de~\Log{G3c}, nous allons
démontrer une généralisation du théorème d'interpolation. Partageons
l'antécédent et le conséquent de chaque séquent en deux parties.
La notation $\maeh{\Gamma_1}{\Gamma_2} \Sequent
\maeh{\Delta_1}{\Delta_2}$ indique que $\Gamma_1$ et $\Delta_1$
constituent la première partie, et $\Gamma_2$ et $\Delta_2$ la seconde.
Le théorème d'interpolation découlera immédiatement du lemme suivant.

\begin{lem}[Lemme de Maehara]\ollabel{lem:maehara}
Si $\Log{G3c} \Proves \maeh{\Gamma_1}{\Gamma_2} \Sequent
\maeh{\Delta_1}{\Delta_2}$, il existe une !!{formula}~$!C$ dont le
!!{language} vérifie $\Lang L(!C) \subseteq \Lang L(\Gamma_1,
\Delta_1) \cap \Lang L(\Gamma_2, \Delta_2)$, telle que
$\Log{G3c} \Proves \Gamma_1 \Sequent \Delta_1, !C$ et
$\Log{G3c} \Proves !C, \Gamma_2 \Sequent \Delta_2$.
\end{lem}

\begin{proof}
Supposons que $\pi \Proves \maeh{\Gamma_1}{\Gamma_2} \Sequent
\maeh{\Delta_1}{\Delta_2}$. Nous raisonnons par récurrence sur
$n=\pheight{\pi}$.

Si $n=0$, le séquent $\maeh{\Gamma_1}{\Gamma_2} \Sequent
\maeh{\Delta_1}{\Delta_2}$ est un axiome. Considérons d'abord un axiome
d'identité : une !!{formula} atomique~$!A$ apparaît à la fois dans
$\Gamma_1, \Gamma_2$ et dans $\Delta_1, \Delta_2$.
Quatre cas se présentent, selon la partie à laquelle appartient
l'occurrence de~$!A$ de chaque côté.
\begin{enumerate}
\item $!A \in \Gamma_1$ et $!A \in \Delta_1$. Nous cherchons une
!!{formula}~$!C$ telle que $\Gamma_1 \Sequent \Delta_1, !C$ et
$!C, \Gamma_2 \Sequent \Delta_2$ soient !!{provable}s. Le premier
séquent est déjà un axiome, quel que soit le choix de~$!C$, puisque
$!A$ figure de chaque côté. Un choix qui garantit uniformément que
le second séquent soit !!{provable} est $!C \ident \lfalse$.
\item $!A \in \Gamma_1$ et $!A \in \Delta_2$. Les séquents
$\Gamma_1 \Sequent \Delta_1, !A$ et $!A, \Gamma_2 \Sequent \Delta_2$
sont alors des axiomes ; nous pouvons donc choisir $!C \ident !A$.
\item $!A \in \Gamma_2$ et $!A \in \Delta_1$. Les séquents
$!A, \Gamma_1 \Sequent \Delta_1$ et
$\Gamma_2 \Sequent \Delta_2, !A$ seraient des axiomes, mais~$!A$
y figure du mauvais côté. Ce choix ne fournit donc pas directement
l'interpolant voulu. À l'aide de $\RightR{\lnot}$ et
$\LeftR{\lnot}$, nous pouvons cependant !!{prove}
$\Gamma_1 \Sequent \Delta_1, \lnot!A$ et
$\lnot !A, \Gamma_2 \Sequent \Delta_2$.
\item $!A \in \Gamma_2$ et $!A \in \Delta_2$. Exercice.
\end{enumerate}
Les cas où $\maeh{\Gamma_1}{\Gamma_2} \Sequent
\maeh{\Delta_1}{\Delta_2}$ est un axiome parce que
$\lfalse \in \Gamma_1$ ou $\lfalse \in \Gamma_2$ sont laissés en exercice.

Si $n>0$, $\pi$ se termine par une inférence logique. Nous distinguons
les cas selon la règle appliquée et la partie à laquelle appartient
la formule principale. L'hypothèse de récurrence s'applique à toutes
les !!{proof}s~$\pi_1$ telles que $\pheight{\pi_1} < n$, quelle que
soit la partition en deux parties de leur séquent final.
Examinons quelques cas.
\begin{enumerate}
\item $\pi$ se termine par $\RightR{\lnot}$, de formule principale
$\lnot !A$. Selon que $\lnot !A$ appartient à la première ou à la
seconde partie, $\pi$ se termine par
\[
    \AxiomC{}
    \RightLabel{$\pi_1$}
    \Deduce$\maeh{!A, \Gamma_1}{\Gamma_2} \fCenter
      \maeh{\Delta_1}{\Delta_2}$
    \RightLabel{\RightR{\lnot}}
    \UnaryInf$\maeh{\Gamma_1}{\Gamma_2} \fCenter
      \maeh{\Delta_1, \lnot !A}{\Delta_2}$
    \DisplayProof
    \]
ou par
\[
    \AxiomC{}
    \RightLabel{$\pi_1$}
    \Deduce$\maeh{\Gamma_1}{!A, \Gamma_2} \fCenter
      \maeh{\Delta_1}{\Delta_2}$
    \RightLabel{\RightR{\lnot}}
    \UnaryInf$\maeh{\Gamma_1}{\Gamma_2} \fCenter
      \maeh{\Delta_1}{\Delta_2, \lnot !A}$
    \DisplayProof
    \]
Lorsque la formule principale~$\lnot !A$ appartient à la première
partie, nous avons placé la formule active~$!A$ dans cette même partie
de la prémisse. Nous sommes libres de faire ce choix. L'hypothèse
de récurrence s'appliquerait aussi en plaçant la formule active dans
l'autre partie, mais l'interpolant obtenu ne pourrait alors,
en général, être repris tel quel pour la conclusion
(exercice : examiner ce placement).

Dans le premier cas, où $\lnot !A$ appartient à la première partie,
l'hypothèse de récurrence donne un interpolant du séquent final
de~$\pi_1$, c'est-à-dire une !!{formula}~$!C_1$ telle que
\begin{align*}
    !A, \Gamma_1 & \Sequent \Delta_1, !C_1 &&\text{et} &
    !C_1, \Gamma_2 & \Sequent \Delta_2
    \intertext{sont !!{provable}s. Il nous faut !!a{formula}~$!C$ telle que les deux séquents suivants soient !!{provable}s :}
    \Gamma_1 & \Sequent \Delta_1, \lnot !A, !C &&\text{et} &
    !C, \Gamma_2 & \Sequent \Delta_2
    \end{align*}
Il suffit de prendre $!C \ident !C_1$. L'hypothèse de récurrence
donne déjà le séquent de droite ; celui de gauche se déduit de la
prémisse correspondante par $\RightR{\lnot}$. Nous savons aussi que
$\Lang L(!C_1) \subseteq \Lang L(!A, \Gamma_1,
\Delta_1) \cap \Lang L(\Gamma_2, \Delta_2)$. Comme
$\Lang L(\lnot !A) = \Lang L(!A)$ et
$\Lang L(!C) = \Lang L(!C_1)$, nous avons
$\Lang L(!C_1) \subseteq \Lang L(\Gamma_1, \Delta_1, \lnot !A) \cap
\Lang L(\Gamma_2, \Delta_2)$.

Dans le second cas, $\lnot !A$ appartient à la seconde partie.
Nous plaçons également la formule active de la prémisse dans cette
partie. L'hypothèse de récurrence donne des preuves des séquents
suivants :
\begin{align*}
    \Gamma_1 & \Sequent \Delta_1, !C_1 &&\text{et} &
    !C_1, !A, \Gamma_2 & \Sequent \Delta_2
    \intertext{sont !!{provable}s. En appliquant de nouveau la règle \RightR{\lnot} au second, nous obtenons des !!{proof}s de}
    \Gamma_1 & \Sequent \Delta_1, !C_1 &&\text{et} &
    !C_1, \Gamma_2 & \Sequent \Delta_2, \lnot !A
    \end{align*}
Ce sont exactement les séquents requis pour que $!C_1$ soit un
interpolant ; nous prenons donc de nouveau $!C \ident !C_1$.
Comme précédemment,
$\Lang L(!C_1) \subseteq \Lang L(\Gamma_1, \Delta_1) \cap
\Lang L(\Gamma_2, \Delta_2, \lnot !A)$.

\item $\pi$ se termine par $\RightR{\lif}$, de formule principale
$!A \lif !B$. Nous avons encore deux cas, selon que $!A \lif !B$
appartient à la première ou à la seconde partie. La !!{proof}~$\pi$
se termine par
\[
    \AxiomC{}
    \RightLabel{$\pi_1$}
    \Deduce$\maeh{!A, \Gamma_1}{\Gamma_2} \fCenter
      \maeh{\Delta_1, !B}{\Delta_2}$
    \RightLabel{\RightR{\lif}}
    \UnaryInf$\maeh{\Gamma_1}{\Gamma_2} \fCenter
      \maeh{\Delta_1, !A \lif !B}{\Delta_2}$
    \DisplayProof
    \]
ou par
\[
    \AxiomC{}
    \RightLabel{$\pi_1$}
    \Deduce$\maeh{\Gamma_1}{!A, \Gamma_2} \fCenter
      \maeh{\Delta_1}{\Delta_2, !B}$
    \RightLabel{\RightR{\lif}}
    \UnaryInf$\maeh{\Gamma_1}{\Gamma_2} \fCenter
      \maeh{\Delta_1}{\Delta_2, !A \lif !B}$
    \DisplayProof
    \]
Nous avons de nouveau placé les formules actives dans la même partie
de la prémisse que la formule principale dans la conclusion.
Dans le premier cas, l'hypothèse de récurrence donne des preuves de
\begin{align*}
    !A, \Gamma_1 & \Sequent \Delta_1, !B, !C_1 &&\text{et} & !C_1, \Gamma_2 & \Sequent \Delta_2
    \intertext{Le séquent de gauche peut servir de prémisse à \RightR{\lif}. Les séquents suivants sont donc !!{provable}s :}
    \Gamma_1 & \Sequent \Delta_1, !A \lif !B, !C_1 &&\text{et} & !C_1, \Gamma_2 & \Sequent \Delta_2
    \end{align*}
Ainsi, $!C_1$ est aussi un interpolant du séquent final de~$\pi$,
et nous pouvons encore prendre $!C \ident !C_1$.

Le second cas se traite de la même façon.

\item $\pi$ se termine par $\LeftR{\lif}$, de formule principale
$!A \lif !B$. Si $!A \lif !B$ appartient à la première partie,
la !!{proof}~$\pi$ se termine par
\[
    \AxiomC{}
    \RightLabel{$\pi_1$}
    \Deduce$\maeh{\Gamma_1}{\Gamma_2} \fCenter
      \maeh{\Delta_1, !A}{\Delta_2}$
    \AxiomC{}
    \RightLabel{$\pi_2$}
    \Deduce$\maeh{!B, \Gamma_1}{\Gamma_2} \fCenter
      \maeh{\Delta_1}{\Delta_2}$
    \RightLabel{\LeftR{\lif}}
    \BinaryInf$\maeh{!A \lif !B, \Gamma_1}{\Gamma_2} \fCenter
      \maeh{\Delta_1}{\Delta_2}$
    \DisplayProof
    \]
L'hypothèse de récurrence donne quatre séquents !!{provable}s :
deux pour l'interpolant $!C_1$ de la prémisse de gauche et deux
pour l'interpolant $!C_2$ de la seconde prémisse :
\begin{align*}
    \Gamma_1 & \Sequent \Delta_1, !A, !C_1 &&\text{(1)} &
    !C_1, \Gamma_2 & \Sequent \Delta_2 && \text{(2)}\\
    !B, \Gamma_1 & \Sequent \Delta_1, !C_2 &&\text{(3)} &
    !C_2, \Gamma_2 & \Sequent \Delta_2 && \text{(4)}
    \end{align*}
Les séquents (1) et (3) peuvent servir de prémisses à
\LeftR{\lif} : par affaiblissement, ajoutons $!C_2$ à droite
dans~(1) et $!C_1$ à droite dans~(3).
\[
    \AxiomC{}
    \Deduce$\Gamma_1 \fCenter \Delta_1, !A, !C_1, !C_2$
    \AxiomC{}
    \Deduce$!B, \Gamma_1 \fCenter \Delta_1, !C_1, !C_2$
    \RightLabel{\LeftR{\lif}}
    \BinaryInf$!A \lif !B, \Gamma_1 \fCenter \Delta_1, !C_1, !C_2$
    \DisplayProof
    \]
En appliquant $\RightR{\lor}$, nous obtenons
\[
    !A \lif !B, \Gamma_1 \fCenter \Delta_1, !C_1 \lor !C_2
    \]
Ceci suggère de choisir $!C_1 \lor !C_2$ comme interpolant.
Le second séquent requis se démontre effectivement à partir
des séquents restants (2) et~(4) :
\[
    \AxiomC{}
    \Deduce$!C_1, \Gamma_2 \fCenter \Delta_2$
    \AxiomC{}
    \Deduce$!C_2, \Gamma_2 \fCenter \Delta_2$
    \RightLabel{\LeftR{\lor}}
    \BinaryInf$!C_1 \lor !C_2, \Gamma_2 \fCenter \Delta_2$
    \DisplayProof
    \]
Il reste à vérifier que $!C_1 \lor !C_2$ appartient au langage voulu.
L'hypothèse de récurrence donne
\begin{align*}
    \Lang L(!C_1) & \subseteq
    \Lang L(\Gamma_1, \Delta_1, !A) \cap \Lang L(\Gamma_2, \Delta_2) &\text{et}\\
    \Lang L(!C_2) & \subseteq
    \Lang L(!B, \Gamma_1, \Delta_1) \cap \Lang L(\Gamma_2, \Delta_2)
    \intertext{Comme}
    \Lang L(!A) & \subseteq \Lang L(!A \lif !B) &\text{et}\\
    \Lang L(!B) & \subseteq \Lang L(!A \lif !B)
    \intertext{nous avons à la fois}
    \Lang L(!C_1) & \subseteq
    \Lang L(\Gamma_1, \Delta_1, !A \lif !B) \cap \Lang L(\Gamma_2, \Delta_2) &\text{et}\\
    \Lang L(!C_2) & \subseteq
    \Lang L(\Gamma_1, \Delta_1, !A \lif !B) \cap \Lang L(\Gamma_2, \Delta_2)
    \intertext{et, par conséquent, $\Lang L (!C_1 \lor !C_2) = {}$}
    \Lang L(!C_1) \cup \Lang L(!C_2) & \subseteq
    \Lang L(\Gamma_1, \Delta_1, !A \lif !B) \cap \Lang L(\Gamma_2, \Delta_2),
    \end{align*}
comme il le fallait.

Lorsque $!A \lif !B$ appartient à la seconde partie, la situation
est différente, mais se traite de façon analogue. L'hypothèse
de récurrence donne encore quatre séquents !!{provable}s :
\begin{align*}
    \Gamma_1 & \Sequent \Delta_1, !C_1 &&\text{(1)} &
    !C_1, \Gamma_2 & \Sequent \Delta_2, !A && \text{(2)}\\
    \Gamma_1 & \Sequent \Delta_1, !C_2 &&\text{(3)} &
    !C_2, !B, \Gamma_2 & \Sequent \Delta_2 && \text{(4)}
    \end{align*}
Les séquents (2) et (4), après quelques affaiblissements, servent
cette fois de prémisses à~\LeftR{\lif} :
\[
    \AxiomC{}
    \Deduce$!C_1, !C_2, \Gamma_2 \fCenter \Delta_2, !A$
    \AxiomC{}
    \Deduce$!B, !C_1, !C_2, \Gamma_2 \fCenter \Delta_2$
    \RightLabel{\LeftR{\lif}}
    \BinaryInf$!A \lif !B, !C_1, !C_2, \Gamma_2 \fCenter \Delta_2$
    \DisplayProof
    \]
Nous voulons de nouveau réunir $!C_1$ et $!C_2$ en une seule formule,
mais cette fois dans l'antécédent, puisque nous traitons la seconde
partie. La règle $\LeftR{\land}$ convient ; nous choisissons
$!C \ident (!C_1 \land !C_2)$ comme interpolant. Les séquents restants
(1) et (3) permettent de démontrer le séquent requis pour l'autre
partie :
\[
    \AxiomC{}
    \Deduce$\Gamma_1 \fCenter \Delta_1, !C_1$
    \AxiomC{}
    \Deduce$\Gamma_1 \fCenter \Delta_1, !C_2$
    \RightLabel{\RightR{\land}}
    \BinaryInf$\Gamma_1 \fCenter \Delta_1, !C_1 \land !C_2$
    \DisplayProof
    \]
Comme précédemment,
$\Lang L(!C_1 \land !C_2) \subseteq \Lang L(\Gamma_1, \Delta_1) \cap
\Lang L(\Gamma_2, \Delta_2, !A \lif !B)$.

\item $\pi$ se termine par $\RightR{\lforall}$, de formule principale
$\lforall[x][!A(x)]$ :
\[
    \AxiomC{}
    \RightLabel{$\pi_1$}
    \Deduce$\maeh{\Gamma_1}{\Gamma_2} \fCenter
    \maeh{\Delta_1, !A(c)}{\Delta_2}$
    \RightLabel{\RightR{\lforall}}
    \UnaryInf$\maeh{\Gamma_1}{\Gamma_2} \fCenter
      \maeh{\Delta_1, \lforall[x][!A(x)]}{\Delta_2}$
    \DisplayProof
    \]
ou
\[
    \AxiomC{}
    \RightLabel{$\pi_1$}
    \Deduce$\maeh{\Gamma_1}{\Gamma_2} \fCenter
      \maeh{\Delta_1}{\Delta_2, !A(c)}$
    \RightLabel{\RightR{\lforall}}
    \UnaryInf$\maeh{\Gamma_1}{\Gamma_2} \fCenter
      \maeh{\Delta_1}{\Delta_2, \lforall[x][!A(x)]}$
    \DisplayProof
    \]
Dans le premier cas, l'hypothèse de récurrence donne des preuves de
$\Gamma_1 \Sequent \Delta_1, !A(c), !C_1$ et
$!C_1, \Gamma_2 \Sequent \Delta_2$. Elle donne aussi
$\Lang L(!C_1) \subseteq \Lang L(\Gamma_1,
\Delta_1, !A(c)) \cap \Lang L(\Gamma_2, \Delta_2)$.
La constante~$c$ ne peut donc pas apparaître dans~$!C_1$ :
sinon, l'inclusion ci-dessus imposerait
$c \in \Lang L(\Gamma_2, \Delta_2)$.
Or $c$ est absente de cette partie de la conclusion par la
condition de constante propre de l'inférence $\RightR{\lforall}$.



Nous pouvons dès lors appliquer $\RightR{\lforall}$ à la première
preuve et obtenir
$\Gamma_1 \Sequent \Delta_1, \lforall[x][!A(x)], !C_1$ :
la constante~$c$ est absente aussi bien du contexte initial que
de l'interpolant ajouté. La condition de vocabulaire
$\Lang L(!C_1) \subseteq \Lang L(\Gamma_1, \Delta_1, \lforall[x][!A(x)])
\cap \Lang L(\Gamma_2, \Delta_2)$ est également satisfaite.
La formule~$!C_1$ est donc un interpolant.

Le second cas est analogue.

\item $\pi$ se termine par $\RightR{\lexists}$, de formule principale
$\lexists[x][!A(x)]$. Nous avons encore deux cas. Nous traitons celui
où $\lexists[x][!A(x)]$ appartient à la première partie et laissons
l'autre en exercice.

Dans le premier cas, $\pi$ se termine par
\[
    \AxiomC{}
    \RightLabel{$\pi_1$}
    \Deduce$\maeh{\Gamma_1}{\Gamma_2} \fCenter
      \maeh{\Delta_1, \lexists[x][!A(x)], !A(t)}{\Delta_2}$
    \RightLabel{\RightR{\lexists}}
    \UnaryInf$\maeh{\Gamma_1}{\Gamma_2} \fCenter
      \maeh{\Delta_1, \lexists[x][!A(x)]}{\Delta_2}$
    \DisplayProof
    \]
L'hypothèse de récurrence donne des preuves de
$\Gamma_1 \Sequent \Delta_1, \lexists[x][!A(x)], !A(t), !C_1$ et
$!C_1, \Gamma_2 \Sequent \Delta_2$. En appliquant
$\RightR{\lexists}$ à la première, nous obtenons
$\Gamma_1 \Sequent \Delta_1, \lexists[x][!A(x)], !C_1$.
Ainsi, $!C_1$ serait encore un interpolant si
$\Lang L(!C_1) \subseteq \Lang L(\Gamma_1, \Delta_1,
\lexists[x][!A(x)]) \cap \Lang L(\Gamma_2, \Delta_2)$.
L'hypothèse de récurrence donne toutefois
$\Lang L(!C_1) \subseteq \Lang L(\Gamma_1, \Delta_1,
\lexists[x][!A(x)], !A(t)) \cap \Lang L(\Gamma_2, \Delta_2)$.

Rappelons qu'il n'y a aucun !!{function} et que les termes
d'instanciation de ce calcul sont clos. Le terme~$t$ est donc
nécessairement une !!{constant}, soit $t \ident b$.
Si $b \notin \Lang L(!C_1)$ ou si
$b \in \Lang L(\Gamma_1, \Delta_1, \lexists[x][!A(x)]) \cap
\Lang L(\Gamma_2, \Delta_2)$, nous pouvons choisir $!C_1$ comme
interpolant. Mais supposons que $b \in \Lang L(!C_1)$ et
$b \notin \Lang L(\Gamma_1, \Delta_1,
\lexists[x][!A(x)]) \cap \Lang L(\Gamma_2, \Delta_2)$.
Remplaçons toutes les occurrences de~$b$ dans~$!C_1$ par une
variable fraîche~$y$, après renommage des variables liées si
nécessaire pour éviter toute capture. Nous pouvons alors écrire
$!C_1 \ident \Subst{!C_1'}{b}{y}$, où $!C_1'$ ne contient pas~$b$.
Nous avons soit
$b \notin \Lang L(\Gamma_1, \Delta_1, \lexists[x][!A(x)])$,
soit $b \notin \Lang L(\Gamma_2, \Delta_2)$.
Dans ces deux cas, nous pouvons prendre respectivement
$!C \ident \lforall[y][!C_1'(y)]$ et
$!C \ident \lexists[y][!C_1'(y)]$.

Dans le premier cas, nous avons
\[
    \AxiomC{}
    \Deduce$\Gamma_1 \fCenter
      \Delta_1, \lexists[x][!A(x)], !A(b), !C_1'(b)$
    \RightLabel{\RightR{\lexists}}
    \UnaryInf$\Gamma_1 \fCenter \Delta_1, \lexists[x][!A(x)], !C_1'(b)$
    \RightLabel{\RightR{\lforall}}
    \UnaryInf$\Gamma_1 \fCenter
      \Delta_1, \lexists[x][!A(x)], \lforall[y][!C_1'(y)]$
    \DisplayProof
    \]
et, par l'hypothèse de récurrence et l'affaiblissement à gauche
par~$\lforall[y][!C_1'(y)]$,
\[
    \AxiomC{}
    \Deduce$!C_1'(b), \lforall[y][!C_1'(y)], \Gamma_2 \fCenter \Delta_2$
    \RightLabel{\LeftR{\lforall}}
    \UnaryInf$\lforall[y][!C_1'(y)], \Gamma_2 \fCenter \Delta_2$
    \DisplayProof
    \]
La condition de constante propre de \RightR{\lforall} dans la
première preuve est satisfaite parce que
$b \notin \Lang L(\Gamma_1, \Delta_1, \lexists[x][!A(x)])$
et parce que $!C_1'$ a été défini de façon à ne pas contenir~$b$.

Dans le second cas, l'hypothèse de récurrence et l'admissibilité
de l'affaiblissement par~$\lexists[y][!C_1'(y)]$ donnent
\[
    \AxiomC{}
    \Deduce$\Gamma_1 \fCenter
      \Delta_1, \lexists[x][!A(x)], !A(b), \lexists[y][!C_1'(y)], !C_1'(b)$
    \RightLabel{\RightR{\lexists}}
    \UnaryInf$\Gamma_1 \fCenter
      \Delta_1, \lexists[x][!A(x)], \lexists[y][!C_1'(y)], !C_1'(b)$
    \RightLabel{\RightR{\lexists}}
    \UnaryInf$\Gamma_1 \fCenter
      \Delta_1, \lexists[x][!A(x)], \lexists[y][!C_1'(y)]$
    \DisplayProof
    \]
et
\[
    \AxiomC{}
    \Deduce$!C_1'(b), \Gamma_2 \fCenter \Delta_2$
    \RightLabel{\LeftR{\lexists}}
    \UnaryInf$\lexists[y][!C_1'(y)], \Gamma_2 \fCenter \Delta_2$
    \DisplayProof
    \]
La condition de constante propre de \LeftR{\lexists} dans la
seconde preuve est satisfaite parce que
$b \notin \Lang L(\Gamma_2, \Delta_2)$
et parce que $!C_1'$ ne contient pas~$b$.
\end{enumerate}
Les autres cas sont analogues ; nous les laissons en exercice.
\end{proof}

\begin{prob}
Supposons que nous disposions d'un connecteur $\lnand$ muni des règles
\[
\Axiom$\Gamma \fCenter \Delta, !A$
\Axiom$\Gamma \fCenter \Delta, !B$
\RightLabel{\LeftR{\lnand}}
\BinaryInf$!A \lnand !B, \Gamma \fCenter \Delta$
\DisplayProof
\quad
\Axiom$!A, !B, \Gamma \fCenter \Delta$
\RightLabel{\RightR{\lnand}}
\UnaryInf$\Gamma \fCenter \Delta, !A \lnand !B$
\DisplayProof
\]
Montrer que le lemme de Maehara reste valable en traitant les cas
de la preuve du \olref{lem:maehara} où $\pi$ se termine par l'une
de ces règles.
\end{prob}

Le lemme de Maehara permet maintenant de démontrer le théorème
d'interpolation de Craig.

\begin{proof}[Preuve du \olref{thm:interpolation}]
Supposons que $!A \Sequent !B$ soit !!{provable}. Appliquons
\olref{lem:maehara} à $\maeh{!A}{\ } \Sequent \maeh{\ }{!B}$
pour obtenir l'interpolant~$!C$. Nous avons
$\Proves !A \Sequent !C$ et $\Proves !C \Sequent !B$.
\end{proof}

Soit $\Gamma$ une théorie, c'est-à-dire un ensemble d'!!{sentence}s,
dans un !!{language}~$\Lang L$ contenant un !!{predicate}~$R$
d'arité $n$. Nous disons que $\Gamma$ \emph{définit implicitement}~$R$
si
\begin{align*}
\Gamma, \Gamma'
& \Entails \lforall[x_1][\dots\lforall[x_n][(R(x_1,\dots,x_n) \liff
R'(x_1, \dots, x_n))]],
\intertext{où $\Gamma'$ s'obtient à partir de $\Gamma$ en remplaçant toutes les occurrences de~$R$ par~$R' \notin \Lang L$. Nous disons que $\Gamma$ \emph{définit explicitement}~$R$ s'il existe une !!{formula}~$!A(x_1, \dots, x_n)$ de ce langage, dont les variables libres sont parmi celles indiquées et qui ne contient pas~$R$, telle que}
\Gamma &\Entails
\lforall[x_1][\dots\lforall[x_n][(R(x_1,\dots,x_n) \liff !A(x_1,
\dots, x_n))]].
\end{align*}
Si $\Gamma$ définit explicitement $R$, elle le définit évidemment
implicitement. La réciproque est moins immédiate, mais elle est vraie :

\begin{lem}[Théorème de définissabilité de Beth]
Si $\Gamma$ définit implicitement~$R$, elle le définit aussi explicitement.
\end{lem}

\begin{proof}
Supposons que $\Gamma$ définisse implicitement~$R$.
Soient $R'$ et $\Gamma'$ comme ci-dessus. La théorie renommée est
dans le langage $\Lang L' = \Lang L \setminus \{R\} \cup \{R'\}$.
Par complétude, l'hypothèse donne une preuve du premier séquent
ci-dessous ; les contextes théoriques désignent ici des sous-ensembles
finis suffisants pour cette preuve.
\begin{align*}
    \Gamma, \Gamma' & 
    \Sequent \lforall[x_1][\dots\lforall[x_n][(R(x_1,\dots,x_n) 
    \liff R'(x_1, \dots, x_n))]]
  \intertext{Soient $c_1$, \dots,~$c_n$ des !!{constant}s distinctes n'apparaissant pas dans~$\Gamma$ et fraîches pour les preuves. Par le lemme d'inversibilité, nous avons}
    \Gamma, \Gamma', R(c_1,\dots,c_n) & \Sequent  R'(c_1, \dots, c_n)
  \intertext{Appliquons le lemme de Maehara à}
    \maeh{\Gamma, R(c_1,\dots,c_n)}{\Gamma'} & \Sequent
      \maeh{\ }{R'(c_1, \dots, c_n)}
  \intertext{pour obtenir un interpolant. En remplaçant dans cet interpolant les constantes introduites par des variables fraîches, sans capture, nous obtenons une !!{formula}~$!A \ident !A(x_1, \dots, x_n)$ telle que}
    \Gamma, R(c_1,\dots,c_n) & \Sequent !A(c_1, \dots, c_n) \\
    !A(c_1, \dots, c_n), \Gamma' & \Sequent R'(c_1, \dots, c_n)
  \intertext{soient !!{provable}s. Le vocabulaire commun des deux parties ne contient ni~$R$ ni le nouveau prédicat. Après abstraction des constantes introduites, $\Lang L(!A) \subseteq \Lang L \setminus \{R\}$. Ainsi, $!A$ ne contient ni le prédicat défini ni son renommage. Remplaçons partout $R'$ par~$R$ dans la preuve du second séquent pour obtenir une preuve de}
    !A(c_1, \dots, c_n), \Gamma & \Sequent R(c_1, \dots, c_n)
  \end{align*}
Appliquons \RightR{\lif} aux deux séquents obtenus, puis la règle
de conjonction à droite pour former l'équivalence, et enfin
\RightR{\lforall} pour fermer universellement les variables.
Les constantes introduites étant fraîches pour le contexte,
nous obtenons une preuve de
\[
    \Gamma \Sequent \lforall[x_1][\dots\lforall[x_n][(R(x_1,\dots,x_n) \liff !A(x_1, \dots, x_n))]].
  \]
Ainsi, $!A(x_1, \dots, x_n)$ définit explicitement $R$ dans~$\Gamma$.
\end{proof}

Un troisième résultat, équivalent à l'interpolation et démontrable
lui aussi à partir du lemme de Maehara, est le suivant.
Si deux théories cohérentes~$\Gamma_1$ et $\Gamma_2$ contiennent une
théorie $\Gamma$ complète pour tous les énoncés du langage commun
$\Lang L(\Gamma_1) \cap \Lang L(\Gamma_2)$,
alors $\Gamma_1 \cup \Gamma_2$ est cohérente.

Rappelons qu'une théorie est complète si, pour tout
!!{sentence}~$!A$ de son langage, soit $\Gamma \Proves !A$, soit
$\Gamma \Proves \lnot !A$. Comme $\Gamma_1$ et $\Gamma_2$ sont
des extensions cohérentes de~$\Gamma$, celle-ci est cohérente.
Il s'ensuit que, pour un !!{sentence}~$!A$ du langage
$\Lang L(\Gamma_1) \cap \Lang L(\Gamma_2)$,
nous ne pouvons avoir à la fois $\Gamma_1 \Entails !A$ et
$\Gamma_2 \Entails \lnot !A$.

En effet, supposons qu'un tel énoncé existe. Si
$\Gamma_1 \Entails !A$, alors $\Gamma \Entails !A$ :
sinon, par complétude de~$\Gamma$, nous aurions
$\Gamma \Entails \lnot !A$ et, puisque $\Gamma \subseteq \Gamma_1$,
$\Gamma_1 \Entails \lnot !A$, ce qui rendrait $\Gamma_1$ incohérente.
Donc $\Gamma \Entails !A$ et, puisque $\Gamma \subseteq \Gamma_2$,
$\Gamma_2 \Entails !A$. Par conséquent,
$\Gamma_2 \Entails/ \lnot !A$, faute de quoi $\Gamma_2$
serait incohérente.

L'absence d'un !!{sentence}~$!A$ du !!{language} commun
$\Lang L(\Gamma_1) \cap \Lang L(\Gamma_2)$ tel que
$\Gamma_1 \Entails !A$ et $\Gamma_2 \Entails \lnot !A$
suffit à la preuve suivante, que nous donnons au moyen du
lemme de Maehara.

\begin{thm}[Théorème de cohérence conjointe de Robinson]
Supposons que $\Gamma_1$ et $\Gamma_2$ soient des théories
cohérentes et qu'il n'existe aucun !!{sentence}~$!A$ du
!!{language} $\Lang L(\Gamma_1) \cap \Lang L(\Gamma_2)$ tel que
$\Gamma_1 \Entails !A$ et $\Gamma_2 \Entails \lnot !A$.
Alors $\Gamma_1 \cup \Gamma_2$ est cohérente.
\end{thm}

\begin{proof}
Supposons au contraire que $\Gamma_1 \cup \Gamma_2$ soit incohérente.
Le séquent $\Gamma_1, \Gamma_2 \Sequent \ $ possède alors une preuve.
Appliquons le lemme de Maehara à
$\maeh{\Gamma_1}{\Gamma_2} \Sequent \maeh{\ }{\ }$ pour obtenir
un !!{sentence}~$!A$ du !!{language}
$\Lang L(\Gamma_1) \cap \Lang L(\Gamma_2)$ tel que
$\Proves \Gamma_1 \Sequent !A$ et
$\Proves !A, \Gamma_2 \Sequent \quad$.
Du second séquent, nous déduisons
$\Proves \Gamma_2 \Sequent \lnot !A$.
Or nous avons supposé qu'aucun énoncé de ce genre n'existe.
\end{proof}

Les contextes des preuves sont finis : dans les arguments précédents,
les théories $\Gamma$, $\Gamma_1$ et $\Gamma_2$ y désignent les
sous-ensembles finis effectivement employés. La compacité permet
d'étendre les résultats aux théories infinies.



















\end{document}
